\documentclass{article}

\usepackage{dibrisunige-report}
\usepackage{codespace}

% Encodings
\usepackage{amsmath,amssymb,gensymb,textcomp}

% Better tables
\usepackage{array,multicol,multirow,siunitx,tabularx}

% Better enum
\usepackage{enumitem}

% Graphics
\usepackage{caption,float}
\usepackage[export]{adjustbox}

% Configurations
\newcounter{memberrowno}
\setcounter{memberrowno}{0}

\ocoursename{Virtual Reality for Robotics - code 104737 - a.y. 2025/2026}
\oreporttype{Report}
\title{Project title}
\oadvisor{Prof. Gianni Vercelli, Prof. Matteo Scaramuzzino}
\reportlayout%

% Custom commands
\newcommand*\mean[1]{\bar{#1}}
% Placeholder for results that are not measured yet: replace every \tbd before submission
\newcommand*\tbd{\textit{TBD}}
% Include a result figure if it exists (keeps the report compiling before the experiments are run)
\newcommand*\resultfig[2][width=0.8\linewidth]{%
  \IfFileExists{#2}{\includegraphics[#1]{#2}}{\fbox{\parbox{0.7\linewidth}{\centering\small figure pending: \texttt{\detokenize{#2}}}}}}

% Title and Metadata
\title{Few-Shot Gaussian Splatting with
Diffusion-Based Data Augmentation}
\author{Mohammadhossein baba}
\date{\today}

\begin{document}

% Cover Page
\coverpage%

\section*{Member list \& Workload}
\begin{center}
  \begin{tabular}{>{\stepcounter{memberrowno}\thememberrowno}llcc}
    \toprule
    \multicolumn{1}{c}{\textbf{No.}} & \textbf{Full name} & \textbf{Student ID} & \textbf{Percentage of work} \\
    \midrule
                                     & Mohammadhossein baba                  & 5919466             & 100\%                       \\
    \bottomrule
  \end{tabular}
\end{center}

\newpage
\tableofcontents
\newpage


% ------------------------------------------------------------------
\section{Introduction}

3D Gaussian Splatting (3DGS)~\cite{kerbl2023gs} is an efficient approach for high-quality novel-view synthesis and
real-time rendering. Its quality depends strongly on the number, coverage and multi-view consistency of the
input images, so reconstruction becomes much harder when only a few real views are available.

This project asks whether synthetic views generated by diffusion models from a \emph{limited} set of real images can
compensate for the missing observations, and when they instead degrade the reconstruction. A dense real capture is
used as a full reference; it is then reduced to a few-shot set (20 real images in the experiments reported here). Synthetic images are
generated \emph{only} from that few-shot set, localized in the real scene, and added to the 3DGS training set at
several synthetic-to-real ratios. All models are evaluated on the same real held-out views.

\paragraph{Contributions.}
(i) A reproducible pipeline on my own capture: HLoc reference map, four open diffusion strategies
(inpainting, outpainting, depth-guided and edge-guided synthesis), localization of the generated images, and a 3DGS
experiment loop. (ii) A full-$N$ baseline and a real-only few-shot baseline evaluated on identical held-out real views,
including views in a region that no condition observed. (iii) A protocol that avoids two silent pitfalls of the stock
3DGS code (Sec.~\ref{sec:protocol}). (iv) Quantitative, cost and visual evidence of when synthetic views help or hurt.


% ------------------------------------------------------------------
\section{Background}

Gaussian Splatting represents a scene by anisotropic 3D Gaussians whose position, scale, opacity, orientation and
colour are optimized from calibrated images; optimization is initialized from the sparse Structure-from-Motion
(SfM) point cloud, and camera poses come from SfM as well~\cite{schonberger2016colmap}. In a few-shot
configuration, coverage is reduced and parts of the geometry are weakly constrained.

Diffusion models~\cite{rombach2022ldm} can synthesize plausible additional views, but a visually convincing image is not
necessarily a geometrically correct observation: small errors in layout, texture or perspective become
contradictory supervision for multi-view optimization. The usefulness of generated views is therefore measured
through their effect on 3DGS, not on image appearance.

The diffusion strategies compared are those recommended in the project brief:
\begin{itemize}[noitemsep]
  \item \textbf{Inpainting}: random regions of a real frame are repainted (occlusion/missing-region filling).
  \item \textbf{Outpainting}: the frame is panned sideways and the newly revealed strip is painted,
        extending the field of view while keeping the layout.
  \item \textbf{Depth-guided synthesis}: ControlNet~\cite{zhang2023controlnet} conditioned on a Depth-Anything-V2~\cite{yang2024depthanythingv2}
        depth map of a slightly panned frame, to approximate a nearby viewpoint.
  \item \textbf{Edge-guided synthesis}: ControlNet conditioned on Canny edges of the panned frame.
\end{itemize}
All four use Stable Diffusion~1.5 backbones in half precision so that they run on a free 16\,GB GPU.


% ------------------------------------------------------------------
\section{Methodology}

\subsection{Data and few-shot construction}
A video of a static indoor scene (82.3\,s, $1920{\times}1080$) was recorded with a slow, continuous camera motion. Frames
were sampled uniformly (about 2.4\,fps), downscaled to 1280\,px width, and the 10\% blurriest (variance of the
Laplacian) were dropped, leaving $180$ real frames (full set, $\geq 100$ as required). A contiguous block of 60 frames
in the middle of the sequence was removed from the incomplete set to simulate a coverage gap. Held-out test views
are every 8th registered real frame outside the gap ($15$ views, close to training frames) plus every 8th frame \emph{inside} the removed gap ($8$ views, in regions no condition has observed); the $23$ views are identical in every condition and are reported overall and separately for the two groups. Since the gap views are held out everywhere, full-$N$ trains on $157$ frames.
The training pool is the remaining $105$ frames. The 20-shot set is 20 frames evenly spread over the pool; the 10- and 5-shot
sets are nested subsets ($5\subset10\subset20$), so differences between sizes are not caused by different frames.

\subsection{Reference map and localization}
A reference SfM map is built from the full real set with HLoc~\cite{sarlin2019hloc}: NetVLAD~\cite{arandjelovic2016netvlad} retrieval,
SuperPoint~\cite{detone2018superpoint} features, SuperGlue~\cite{sarlin2020superglue} matching and COLMAP triangulation with a single shared camera.
Each synthetic image is resized to the map resolution and localized against this map with the map's intrinsics.
The number of PnP inliers is recorded and optionally used to filter low-quality localizations. Synthetic
images whose source frame is a test frame are excluded (leakage guard; by construction sources are only
few-shot training frames). Filenames encode the strategy, the source real frame and the variant
(\texttt{<strategy>\_frame\_XXXX\_<variant>.jpg}), linking each synthetic view to its approximate real view.

\subsection{Conditions and training}
Conditions are: the full-$N$ baseline (all $157$ real training frames); the real-only few-shot baseline (0\%);
and few-shot plus $S=\mathrm{round}(N\cdot r)$ synthetic images, with $r\in\{25,50,100\}\%$ for each of the four strategies. Synthetic images are added only to the training side. The same
graphdeco 3DGS implementation and hyper-parameters are used for every condition
(15\,000 iterations, half resolution, \texttt{--init-points visible}). Experiments ran on a free Google Colab GPU (NVIDIA T4, about 15\,GB).

\subsection{Evaluation protocol}
\label{sec:protocol}
Metrics on the held-out real views are PSNR, SSIM and LPIPS~\cite{zhang2018lpips}; training time and peak GPU memory are
logged per condition. Two protocol pitfalls of the stock code were avoided:
\begin{enumerate}[noitemsep]
  \item In the stock 3DGS code the evaluation split is every 8th image by sorted name and the provided
        \texttt{test.txt} is ignored, which can place synthetic images in the test set. The code was patched
        to read \texttt{test.txt}, and the run aborts unless the number of rendered test views equals its length.
  \item The initial point cloud is restricted to points observed by at least two \emph{selected real training}
        images, instead of copying the full map's points into every few-shot scene (which leaks geometry from
        frames outside the few-shot set).
\end{enumerate}


% ------------------------------------------------------------------
\section{Experimental Setup}

\begin{table}[H]
\centering
\caption{Experiment grid (20-shot setting, 14 conditions).}
\begin{tabular}{ll}
\toprule
Factor & Values \\
\midrule
Few-shot size $N$ & 20 (5- and 10-shot sets are defined in the code but were not run) \\
Strategy & inpainting, outpainting, depth-guided, edge-guided \\
Synthetic ratio & 0\%, 25\%, 50\%, 100\% \\
Baselines & full-$N$ (157 real), few-shot real-only \\
Test views & 23 real held-out views (15 near-train, 8 in-gap), identical everywhere \\
\bottomrule
\end{tabular}
\end{table}

All runs share the GS configuration. Localization statistics (success rate, PnP inliers) are collected per
strategy to relate the geometric consistency of generated views to the final reconstruction quality.


% ------------------------------------------------------------------
\section{Results}

\subsection{Quantitative results}

Table~\ref{tab:res} lists all conditions. ``Near'' and ``gap'' are PSNR on the 15 test views next to training frames and the 8
test views inside the removed block. Depth-guided synthesis at 25\% and 100\% diverged during training (loss above 0.4,
GPU memory exhausted at iteration 5\,300 and 6\,800) and produced no model; they are shown as ``fail''.

\begin{table}[H]
\centering
\caption{20-shot GS results on the 23 held-out real views (15\,000 iterations, half resolution, Colab T4). Best value of each metric in bold.}
\label{tab:res}
\small
\resizebox{\linewidth}{!}{%
\begin{tabular}{lccccccc}
\toprule
Condition & Train img. & Time [min] & Mem. [MB] & PSNR $\uparrow$ & SSIM $\uparrow$ & LPIPS $\downarrow$ & PSNR near / gap \\
\midrule
full-$N$           & 157 & 19.0 & 5001 & \textbf{36.74} & \textbf{0.969} & \textbf{0.167} & 37.36 / 35.59 \\
20-shot real only  &  20 & 17.0 & 3941 & 30.47 & 0.936 & 0.209 & 34.61 / 22.73 \\
\midrule
inpaint 25\%       &  25 & 16.4 & 3337 & 30.23 & 0.934 & 0.217 & 34.22 / 22.75 \\
inpaint 50\%       &  30 & 16.1 & 3433 & 29.28 & 0.930 & 0.226 & 33.22 / 21.90 \\
inpaint 100\%      &  40 & 16.1 & 3361 & 28.09 & 0.926 & 0.242 & 31.74 / 21.23 \\
\midrule
outpaint 25\%      &  25 & 17.1 & 3375 & 28.59 & 0.922 & 0.238 & 32.15 / 21.91 \\
outpaint 50\%      &  30 & 17.3 & 3873 & 25.96 & 0.913 & 0.258 & 28.71 / 20.82 \\
outpaint 100\%     &  40 & 17.2 & 3783 & 24.61 & 0.895 & 0.292 & 27.72 / 18.77 \\
\midrule
depth 25\%         &  25 & \multicolumn{6}{c}{fail (diverged, out of GPU memory)} \\
depth 50\%         &  30 & 36.9 & 9083 & 18.42 & 0.849 & 0.426 & 20.08 / 15.31 \\
depth 100\%        &  40 & \multicolumn{6}{c}{fail (diverged, out of GPU memory)} \\
\midrule
edge 25\%          &  25 & 17.5 & 2953 & 25.54 & 0.902 & 0.270 & 28.53 / 19.92 \\
edge 50\%          &  30 & 18.1 & 3259 & 19.68 & 0.817 & 0.363 & 22.65 / 14.11 \\
edge 100\%         &  40 & 16.5 & 2909 & 15.78 & 0.788 & 0.440 & 17.03 / 13.44 \\
\bottomrule
\end{tabular}}
\end{table}

\subsection{Metric plots}
\begin{figure}[H]
\centering
\resultfig[width=0.48\linewidth]{results/plot_PSNR_few20.png}\hfill
\resultfig[width=0.48\linewidth]{results/plot_LPIPS_few20.png}
\caption{PSNR and LPIPS versus synthetic ratio per strategy (20-shot). Dotted line: full-$N$ baseline. Depth-guided has no points at 25\% and 100\% because those runs diverged.}
\end{figure}

\subsection{Main observations}
\begin{itemize}[noitemsep]
  \item \textbf{Headroom.} Full-$N$ reaches 36.7\,dB, 6.3\,dB above the 20-shot real-only baseline (30.5\,dB). The gap is
        much larger on the views inside the removed block (35.6 versus 22.7\,dB), where the few-shot model has no
        nearby observation. This is the room that synthetic views could in principle fill.
  \item \textbf{No strategy beat the real-only baseline.} Every synthetic condition has lower PSNR, SSIM and higher
        LPIPS than 20-shot real only. The only near tie is inpainting at 25\% (30.23 versus 30.47\,dB; 22.75 versus
        22.73\,dB in the gap), a difference that I cannot separate from noise with a single run.
  \item \textbf{More synthetic images hurt more.} For all strategies that trained, PSNR decreases monotonically with the
        ratio (inpaint: 30.2, 29.3, 28.1\,dB; outpaint: 28.6, 26.0, 24.6\,dB; edge: 25.5, 19.7, 15.8\,dB).
  \item \textbf{Strategy ranking.} Inpainting $>$ outpainting $>$ edge-guided $>$ depth-guided. The ranking follows
        how much the generated image departs from the real source view: inpainting keeps most pixels of a real frame,
        outpainting adds a new strip, and the two ControlNet strategies regenerate the whole image.
  \item \textbf{Training fit versus test quality.} With synthetic data the training PSNR stays high at low ratios
        (inpaint 25\%: 40.2\,dB, outpaint 25\%: 39.5\,dB, against 42.1\,dB for real only) while the test PSNR falls: the
        model fits the generated views but they are not consistent with the real geometry. At high ratios even the
        training fit degrades (edge 100\%: 18.2\,dB, depth 50\%: 25.5\,dB), showing contradictory supervision.
  \item \textbf{Cost.} Training time is about 16--18\,min and 3--4\,GB peak memory for every converged few-shot run,
        so synthetic images add little cost (the generation time of the diffusion models is not included). Full-$N$
        needs 19\,min and 5.0\,GB. Depth 50\% was the exception (36.9\,min, 9.1\,GB) presumably because the number of Gaussians grew
        abnormally while the optimization struggled (not logged); two depth runs did not finish.
\end{itemize}

\subsection{Localization of synthetic views}
All 40 generated images of every strategy (20 sources, 2 variants each) were localized in the real map
(success rate 100\% for the four strategies). This number is not discriminative: feature matching returns a
pose even for images that are geometrically wrong, so success is a necessary but not a sufficient condition for a
useful training view. Per-image PnP inlier counts are written by the injection script, but I did not use them to
filter images in the reported runs; an inlier-based selection is the natural next experiment.


% ------------------------------------------------------------------
\section{Discussion}

In this scene and with this protocol, diffusion-generated views did not help few-shot 3DGS; they were neutral at best
(inpainting at 25\%) and increasingly harmful with more images. There is therefore no optimal ratio above 25\%, and
ratios below 25\% (not tested) are the only region where a small benefit could still exist.

The results suggest that the benefit depends on geometric consistency, not on perceptual quality. A 3DGS optimization treats each training
image as an exact observation from its estimated pose. A generated image that looks plausible but differs in
layout or texture from what the pose implies adds contradictory constraints; the optimizer either fits it with
view-dependent floaters, which lowers test quality, or fails to converge (depth-guided runs). The two ControlNet
strategies regenerate the entire image from a depth or edge map of a panned frame; the generated colours, materials and objects
are invented independently in each image and differ from the real scene (Fig.~\ref{fig:synth}), which explains why they are the
least consistent. Inpainting leaves most pixels
equal to the real frame, so its views stay close to real observations and cost the least.

The in-gap views show that the missing coverage is exactly where the augmentation was meant to help; yet the best synthetic
condition did not improve them (22.75 versus 22.73\,dB), because the generated views are derived from frames that
do not observe that region.

Only the 20-shot setting was run; whether augmentation helps more with 5 shots, where real coverage is much lower,
remains untested.


% ------------------------------------------------------------------
\section{Qualitative observations}

\subsection{What the generators produce}
Figure~\ref{fig:synth} shows three source frames with one synthetic image from each strategy. The strategies depart from the real
frame by very different amounts. \textbf{Inpainting} keeps the wall, table and radiator but invents new objects on the table
(a laptop, a notebook, a chair); \textbf{outpainting} keeps the centre of the frame and adds books and furniture at the
border. Both stay recognizably the same scene. The \textbf{depth-} and \textbf{edge-guided} images keep only a coarse
layout: the white table becomes dark grey or brown wood, a laptop and blinds appear, and the colours and materials differ
from the real scene in every example. These images are plausible photographs of \emph{some} desk, but not observations of \emph{this}
desk, so they cannot be consistent with the real views from any camera pose.

\begin{figure}[H]
\centering
\resultfig[width=\linewidth]{results/synthetic_examples.png}
\caption{Real source frames (left) and one synthetic image per strategy generated from each of them.}
\label{fig:synth}
\end{figure}

\subsection{Effect on the rendered views}
Figure~\ref{fig:panel} shows held-out views rendered by each model at 50\% synthetic images, with per-view PSNR. Full-$N$ is close to the
ground truth. The real-only model is sharp where training frames are near (first row, 35.6\,dB) but shows floaters and
smearing in the view that lies inside the removed block (second row, 20.7\,dB). Inpainting at 50\% is visually close to the real-only model
(33.9 versus 35.6\,dB on the first row) with slightly more smearing; outpainting adds visible haze and floaters at the image borders, where its
generated content lives. The edge-guided model shows glass-like streaks and fragments and the depth-guided model produces washed-out, blurred
surfaces in which the scene structure is lost (15--22\,dB). This matches the metrics: the more the generated content
differs from the real scene, the stronger the artifacts.

\begin{figure}[H]
\centering
\resultfig[width=\linewidth]{results/panel_few20_ratio50.png}
\caption{Held-out views (rows). Columns: ground truth, full-$N$, 20-shot real only, and each strategy at 50\%; per-view PSNR in the corner.}
\label{fig:panel}
\end{figure}


% ------------------------------------------------------------------
\section{Limitations}

\begin{itemize}[noitemsep]
  \item One static indoor scene and one capture; no repetitions over seeds, so small metric differences may be noise.
  \item Free-tier compute: reduced resolution and iterations relative to the original 3DGS setting. Only the 20-shot
        set and ratios up to 100\% were run; 5/10-shot and 200\% were left out for time.
  \item Stable Diffusion 1.5 backbones generate at about 768\,px and are upsampled to the map resolution; stronger
        backbones or video models could behave differently.
  \item Synthetic views are localized with feature matching; images whose content was changed (inpainting) may be
        localized less reliably, and a mislocalized image acts as noise, not as extra information.
  \item Radial distortion of the estimated camera is ignored, as 3DGS requires pinhole cameras.
\end{itemize}


% ------------------------------------------------------------------
\section{Conclusion}

Diffusion-based augmentation was tested on few-shot 3DGS with four strategies on my own capture. Full-$N$ reached
36.7\,dB and 20-shot real-only 30.5\,dB, but no synthetic condition improved on real-only. Inpainting at 25\% was neutral,
every strategy lost quality as the ratio grew, and depth- and edge-guided synthesis were the most harmful (two depth runs
diverged).  Future work: filtering
synthetic views by PnP inliers or reprojection error, ratios below 25\%, 5/10-shot sets, multi-view-consistent generators
(video or multi-view diffusion), and repetitions over seeds and scenes.

\section*{Code availability}
Repository: \texttt{github.com/mohammadhoseinbaba/Few\_Shot\_Gaussian\_Splatting\_with\_Diffusion\_Based\_Data\_Augmentation}.

\bibliographystyle{plain}
\bibliography{refs}
\end{document}
